% The only source of this talk. `make TALK=breeden-litzenberger` builds
% slides, handout, transparencies, article, pdfpc notes and a printable script.
%
% Conventions (see README.md):
%   - Text outside frames appears only in the article.
%   - Every frame carries a \note{...} for the presenter.
%   - Numbers come from data/numbers.tex (generated by derive.py), never typed.
\documentclass[
  beameroptions={aspectratio=169,ignorenonframetext,11pt},
  articleoptions={11pt,a4paper},
]{beamerswitch}
\usepackage{yjtalk}
% Generated by derive.py from data/raw.json. Do not edit.
\newcommand{\blSzero}{100}
\newcommand{\blRate}{5\%}
\newcommand{\blSigma}{20\%}
\newcommand{\blT}{1}
\newcommand{\blKzero}{100}
\newcommand{\blErT}{1.0513}
\newcommand{\blCzero}{10.4506}
\newcommand{\blPzero}{0.019724}
\newcommand{\blCpp}{0.018762}
\newcommand{\blEps}{0.01}
\newcommand{\blDeltaMin}{0.5}
\newcommand{\blEstMin}{0.019723}
\newcommand{\blNoiseMin}{0.1682}
\newcommand{\blNoiseRatio}{9}
\newcommand{\blStrikes}{201}


\title{Reading a probability out of option prices}
\subtitle{The Breeden--Litzenberger formula, and why it breaks on real quotes}
\author{Teo Yu Jie}
\institute{Example talk for the beamerswitch template}
\date{29 September 2026}

\pgfplotstableread[col sep=comma]{data/curves.csv}\curves
\pgfplotstableread[col sep=comma]{data/butterflies.csv}\butterflies

\begin{document}

\begin{frame}[plain]
  \titlepage
  \note{Open with the question, not the formula: can a list of option prices
    tell you what the market thinks the stock will be worth in a year? You
    will answer yes, show why in two derivatives, then show where it breaks.
    About 12 minutes in total.}
\end{frame}

\mode<article>{\maketitle}

\section{The question}

An option price is a bet on where a stock ends up. Many bets at many strikes
should, together, say how likely each ending price is. This talk makes that
precise and shows the one place the precise version fails in practice.

\begin{frame}{A call option pays the part of the price above the strike}
  \begin{columns}[c]
    \column{0.5\textwidth}
      \begin{itemize}
        \item Pays $(S_T - K)_+ = \max\{S_T - K,\,0\}$ at time $T$.
        \item<2-> Price today: $C(K) = e^{-rT}\,\mathbb{E}\big[(S_T - K)_+\big]$.
        \item<3-> One price per strike $K$. Many strikes give a curve.
      \end{itemize}
    \column{0.45\textwidth}
      \begin{tikzpicture}
        \begin{axis}[yj, width=6.5cm, height=4.6cm,
            xlabel={stock price at expiry $S_T$}, ylabel={payoff},
            xmin=60, xmax=140, ymin=0, ymax=42, xtick={100}, xticklabels={$K$}, ytick=\empty]
          \addplot[yjForeground, domain=60:140, samples=3] {max(x-100,0)};
          \node[anchor=south east, font=\footnotesize, text=yjSelected] at (axis cs:100,1) {kink};
          \fill[yjSelected] (axis cs:100,0) circle[radius=2pt];
        \end{axis}
      \end{tikzpicture}
  \end{columns}
  \note{Stress the kink: the payoff is flat, then rises with slope one. Every
    result in the talk comes from that corner. Point at the red dot.
    Transition: if the payoff has a corner at $K$, what does the price do as
    you slide $K$?}
\end{frame}

\section{Two derivatives}

Differentiate the price under the expectation. The payoff's first derivative in
$K$ is minus a step, its second is a Dirac delta at $S_T = K$, and the delta
picks the density out of the integral.

\begin{frame}{Differentiating twice in the strike isolates the density}
  \begin{align*}
    C(K) &= e^{-rT}\!\int_0^\infty (s-K)_+\,p(s)\,ds \\
    \uncover<2->{C'(K) &= -e^{-rT}\!\int_K^\infty p(s)\,ds
      && \text{\small\color{yjSecondary} slope: minus a step}\\}
    \uncover<3->{C''(K) &= e^{-rT}\,p(K)
      && \text{\small\color{yjSecondary} curvature: a delta at $s=K$}}
  \end{align*}
  \uncover<4->{\centering\yjkey{$p(K) = e^{rT}\,C''(K)$}: the density is the
    curvature of the call-price curve.\par}
  \note{Reveal one line per click. On line two, say it in words: raising the
    strike by a dollar costs you a dollar in every state where the option pays.
    On line three, the step becomes a spike at $K$. Pause on the red line; it is
    the whole talk. Transition: does real curvature look like a density?}
\end{frame}

\begin{frame}{The call-price curve bends most where the stock is most likely to end}
  \centering
  \begin{tikzpicture}
    \begin{axis}[yj, width=12cm, height=6cm, xmin=60, xmax=160,
        xlabel={strike $K$}, ylabel={call price $C(K)$}, ymin=0]
      \addplot[yjForeground] table[x=K, y=call] {\curves};
      \draw[yjSecondary, dashed] (axis cs:\blKzero,0) -- (axis cs:\blKzero,\blCzero);
      \fill[yjSelected] (axis cs:\blKzero,\blCzero) circle[radius=2.5pt];
      \node[anchor=south west, font=\footnotesize, text=yjSelected]
        at (axis cs:\blKzero,\blCzero) {$C(\blKzero) = \blCzero$};
    \end{axis}
  \end{tikzpicture}
  \yjsource{Synthetic Black--Scholes prices, $S_0=\blSzero$, $r=\blRate$,
    $\sigma=\blSigma$, $T=\blT$; \blStrikes\ strikes. Not market data.}
  \note{Point at the region around $K=\blKzero$: the curve is most convex near
    today's price. Far out, it is nearly straight. Say plainly that these are
    model prices, not quotes; that matters on slide 7. Transition: turn
    curvature into a number.}
\end{frame}

\begin{frame}{Its second difference lands on the density at every strike}
  \centering
  \begin{tikzpicture}
    \begin{axis}[yj, width=12cm, height=6cm, xmin=60, xmax=160,
        xlabel={strike $K$}, ylabel={density $p(K)$}, ymin=0,
        yticklabel style={/pgf/number format/fixed, /pgf/number format/precision=3}]
      \addplot[yjBorder, line width=5pt] table[x=K, y=density] {\curves};
      \addplot[yjSelected, line width=1pt] table[x=K, y=fd] {\curves};
      \node[anchor=west, font=\footnotesize, text=yjSecondary] at (axis cs:118,0.017) {closed form (grey)};
      \node[anchor=west, font=\footnotesize, text=yjSelected] at (axis cs:118,0.0145) {$e^{rT}\,\Delta^2 C / h^2$ (red)};
    \end{axis}
  \end{tikzpicture}
  \yjsource{Second difference with $h=0.5$ on the same synthetic prices.}
  \note{The thick grey band is the lognormal density; the thin red line is
    computed only from call prices. They coincide. At $K=\blKzero$ the density
    is \blPzero. Likely question: why lognormal? Answer: the prices are
    Black--Scholes, so the implied density must be lognormal; this is a check,
    not a discovery. Transition: can you trade a second difference?}
\end{frame}

\section{The butterfly}

A second difference of prices is itself a price: of a portfolio long one call
at $K-\Delta$, short two at $K$ and long one at $K+\Delta$. Its payoff is a
tent of area $\Delta^2$, so its price divided by $\Delta^2$ is the discounted
probability of landing near $K$.

\begin{frame}{A butterfly spread is a tradeable second difference}
  \[ C(K{-}\Delta) - 2\,C(K) + C(K{+}\Delta) \approx \Delta^2\, C''(K) \]
  \begin{columns}[c]
    \column{0.48\textwidth}
      \begin{itemize}
        \item Long one call at $K{\pm}\Delta$, short two at $K$.
        \item<2-> Payoff: a tent of height $\Delta$, area $\Delta^2$.
        \item<3-> Price $\div\,\Delta^2 \times e^{rT}$ estimates $p(K)$.
      \end{itemize}
    \column{0.48\textwidth}
      \begin{tikzpicture}
        \begin{axis}[yj, width=6.5cm, height=4.6cm, xmin=85, xmax=115, ymin=0, ymax=7,
            xtick={95,100,105}, xticklabels={$K{-}\Delta$,$K$,$K{+}\Delta$}, ytick=\empty,
            xlabel={$S_T$}, ylabel={payoff}]
          \addplot[yjMark, fill=yjMark, fill opacity=0.15] coordinates {(85,0) (95,0) (100,5) (105,0) (115,0)};
          \node[font=\footnotesize, text=yjMark] at (axis cs:100,6) {height $\Delta$};
        \end{axis}
      \end{tikzpicture}
  \end{columns}
  \note{Draw the tent in the air with your hand. The point: nothing here needs
    a model; it is three quotes and arithmetic. Transition: shrink $\Delta$ and
    watch the estimate converge.}
\end{frame}

\begin{frame}{Narrower butterflies converge on the density}
  \centering
  \pgfplotstabletypeset[
    columns={delta,price,estimate,bias},
    columns/delta/.style={column name={half-width $\Delta$}, fixed, precision=1, zerofill},
    columns/price/.style={column name={butterfly price}, fixed, precision=6, zerofill},
    columns/estimate/.style={column name={$e^{rT}\,\text{price}/\Delta^2$}, fixed, precision=6, zerofill},
    columns/bias/.style={column name={error vs.\ $p(\blKzero)$}, fixed, precision=6, zerofill},
    every head row/.style={before row=\toprule, after row=\midrule},
    every last row/.style={after row=\bottomrule},
  ]{\butterflies}
  \par\bigskip
  Closed form: $p(\blKzero) = \yjkey{\blPzero}$. At $\Delta=\blDeltaMin$ the
  butterfly gives \blEstMin.
  \yjsource{Same synthetic prices; $K=\blKzero$.}
  \note{Read one row, not all four: at half-width \blDeltaMin\ the estimate
    matches to the sixth decimal. The error column shrinks roughly with
    $\Delta^2$. Transition: so just use the narrowest butterfly? That is the
    trap.}
\end{frame}

\section{Where it breaks}

On real quotes each price carries an error of about the tick size $\varepsilon$.
The butterfly's error is then up to $4\varepsilon$, and dividing by $\Delta^2$
amplifies it. Bias falls as $\Delta$ shrinks; noise grows.

\begin{frame}{On real quotes, shrinking $\Delta$ trades bias for noise}
  \centering
  \begin{tikzpicture}
    \begin{axis}[yj, width=9.5cm, height=5.4cm, xmode=log, ymode=log,
        xlabel={half-width $\Delta$}, ylabel={error in $p(K)$},
        xtick={0.5,1,2,5}, xticklabels={0.5,1,2,5}, x dir=reverse]
      \addplot[yjSecondary, mark=*] table[x=delta, y=bias] {\butterflies};
      \addplot[yjSelected, mark=*] table[x=delta, y=noise] {\butterflies};
      \node[anchor=west, font=\footnotesize, text=yjSecondary] at (axis cs:0.5,0.000001) {bias};
      \node[anchor=west, font=\footnotesize, text=yjSelected] at (axis cs:0.5,0.1682) {noise bound, $\varepsilon=\blEps$};
    \end{axis}
  \end{tikzpicture}
  \par At $\Delta=\blDeltaMin$ the noise bound is \blNoiseMin: \yjkey{about
  \blNoiseRatio\ times the density itself}.
  \yjsource{Bias from the synthetic prices; noise bound $e^{rT}\,4\varepsilon/\Delta^2$
    with an assumed quote error $\varepsilon=\blEps$.}
  \note{This is the turn. Everything so far was exact because the prices were
    exact. With one-cent quote errors, the narrowest butterfly is useless.
    Say that $\varepsilon$ is an assumption for illustration, not measured.
    Likely question: what do practitioners do? Answer: fit a smooth curve to
    implied volatilities first, then differentiate the fit. Transition: so
    what should you remember?}
\end{frame}

\begin{frame}{What to remember}
  \begin{enumerate}
    \item The density is the curvature: $p(K) = e^{rT}\,C''(K)$.
    \item<2-> A butterfly spread is that curvature, as a trade.
    \item<3-> Differentiating noisy prices twice amplifies noise by $1/\Delta^2$:
      smooth before you differentiate.
  \end{enumerate}
  \vfill
  {\scriptsize\color{yjSecondary}
    Breeden \& Litzenberger (1978), \emph{Prices of state-contingent claims
    implicit in option prices}, J.\ Business 51(4).\quad
    Martin (2018), \emph{Options and the gamma knife}, JPM.\par}
  \note{Say the three lines slowly, one per click, then stop talking. Closing
    line: ``Option prices already contain a probability distribution; the
    hard part is not the calculus, it is the noise.'' Take questions.}
\end{frame}

\end{document}
